\subsection{Sprint 4}

Na Sprint 4, as atividades planejadas pelas quais fiquei responsável incluíam o desenvolvimento e a integração da tela do diário da mãe, além do fornecimento de suporte a colegas com menos experiência. Essas atividades foram definidas com o objetivo de avançar na finalização de uma das funcionalidades principais do projeto, garantindo que o produto estivesse alinhado às expectativas das usuárias e aos requisitos estabelecidos antes da entrega final. O trabalho dessa sprint foi intensivo, resultando na finalização das tarefas de desenvolvimento a integração da tela de diário antecipadamente. Isso abriu espaço para eu pudesse trabalhar em outras tarefas que se encontravam somente no \textit{backlog} do projeto, como a criação e documentação de testes unitários no front-end e a refatoração da tela de informações, que sofreu alterações de última hora e precisava ter sua estilização ajustada.

Apesar das conquistas, alguns desafios marcaram essa sprint, especialmente no que diz respeito à procrastinação com relação aos testes. Embora os testes unitários tenham sido desenvolvidos e documentados, a falta de prioridade dada a eles desde o início do projeto resultou na impossibilidade de cobrir todos os requisitos esperados. Além disso, o tempo limitado impediu a implementação de testes de integração, que são essenciais para validar a comunicação entre diferentes partes do sistema. Essa experiência ressaltou a importância de tratar os testes como parte fundamental do ciclo de desenvolvimento, abordando-os de maneira contínua e consistente desde o início do projeto para evitar acúmulos ou lacunas próximas ao prazo final.

Desse modo, o principal aprendizado desta sprint foi a necessidade de adotar uma abordagem proativa em relação aos testes desde o início do projeto. Incorporar a preocupação com cobertura e integração ao longo de todo o desenvolvimento é essencial para garantir um produto robusto e confiável.

Com o fim da Sprint 4, as atenções da equipe agora se voltam apenas para a apresentação final dos projetos da \ac{ages} e a elaboração do relatório final do projeto que sintetiza os esforços, aprendizados e resultados obtidos ao longo do desenvolvimento. Assim,esta sprint, apesar dos obstáculos, demonstrou a capacidade do time de entregar valor e de aprender com os desafios enfrentados, contribuindo para a consolidação de um produto de qualidade.